% \RequirePackage[iso]{isodateo}
\RequirePackage{datetime}
\renewcommand{\dateseparator}{-}
\newdateformat{iso}{\THEYEAR-\twodigit{\THEMONTH}-\twodigit{\THEDAY}}
\newtimeformat{nummeric}{\twodigit{\THEHOUR}\twodigit{\THEMINUTE}}
\settimeformat{nummeric}

\newcommand{\VarMkatalogVersion}{2013}
\newcommand{\VarAassVersion}{1.0}
\newcommand{\DocVersion}{\pdfdate}
% \renewcommand{\DocVersion}{2013-01-09}
\newcommand{\DocVersionTyp}{IN ARBEIT}
% % \renewcommand{\DocVersionTyp}{Vorabversion}
% % \renewcommand{\DocVersionTyp}{Beta-Version}
% \renewcommand{\DocVersionTyp}{FREIGEGEBEN}
\newcommand{\DocVersionInfo}{IN ARBEIT}
% % \renewcommand{\DocVersionInfo}{Vorabversion -- F.\,d. internen Gebrauch.} 
% % \renewcommand{\DocVersionInfo}{Beta-Version -- Zur Ansicht}
% \renewcommand{\DocVersionInfo}{FREIGEGEBEN}
\hypersetup{
% draft,
unicode=true,
colorlinks=false,
linkcolor=darkBlue,
citecolor=darkBlue,
filecolor=darkBlue,
pagecolor=darkBlue,
urlcolor=darkBlue,
% Rahmen um Links:
pdfborderstyle={/S/U/W 0},  % Unterstreichung
% pdfborderstyle={/S/U/W 0.5},  % Unterstreichung
pdfborder={0 0 0},          % Dicke
% linkbordercolor=DefaultB,     % Farbe
% urlbordercolor=DefaultB,     % Farbe
pdftitle={Arbeits- und Ausbildungsstandards für den Sanitätsdienst, Wartungsnotiz für AASS {\VarAassVersion} und Massnahmenkatalog {\VarMkatalogVersion}},
% Sets the document information Title field pdfauthor=ARGE AASS, % Sets the document information Author field
pdfsubject={AASS Wartungsnotiz für AASS {\VarAassVersion} und Massnahmenkatalog {\VarMkatalogVersion}}, % Sets the document information Subject field
pdfcreator=ARGE AASS, % Sets the document information Creator field
pdfproducer=ARGE AASS, % Sets the document information Producer field
pdfkeywords=ARGE AASS, % Sets the document information Keywords field
% pdfpagelayout=TwoPageRight,
}
% By default, URLs are printed using mono-spaced fonts. If you
% don't like it and you want them to be printed with the same style 
% of the rest of the text, you can use this:
\urlstyle{same}
% Staus Möglichkeiten:
% Final 				Endversion
% DraftLight			+ Bilder, - Kommentare
% DraftMedium		- Bilder, - Kommentare
% DraftFull			- Bilder, + Kommenatre
% DraftSuper			- Bilder, + Kommenatare, + draft=yes
\newcommand{\DocStatus}{DraftFull}
\newcommand{\DocGitCommit}{No SHA1!}
% \renewcommand{\DocGitCommit}{(0) \input{../tmp/git-lastcommitsha1.tmp}}
\newcommand{\DocGitCommitFormatted}{\ttfamily\DocGitCommit}
\newcommand{\DocGitCommitTiny}{\tiny\DocGitCommitFormatted}
\newcommand{\DocGitCommitScriptsize}{\scriptsize\DocGitCommitFormatted}
\newcommand{\DocGitCommitFootnotesize}{\footnotesize\DocGitCommitFormatted}